%%%%%%%%%%%%%%%%%%%%%%%%%%%%%%%%%%%%%%%%%%%%%%%%%%%%%%%%%%%%%%%%%%%%%%%%%%%%%%%
%% Resumo --- ABNT NBR 6028:2021
%%   parágrafo único; voz ativa; sem citações;
%%   palavras-chave separadas por ponto e vírgula, iniciais minúsculas,
%%   finalizadas por ponto.
%%
%% Ao alterar a resposta ao problema de pesquisa, alterar também as Conclusões e o
%% abstract.
%%%%%%%%%%%%%%%%%%%%%%%%%%%%%%%%%%%%%%%%%%%%%%%%%%%%%%%%%%%%%%%%%%%%%%%%%%%%%%%

\begin{resumo}
A análise de vibração permite decidir a manutenção de motores pelo estado medido da máquina,
mas o custo da instrumentação restringe seu uso às máquinas críticas e deixa a maior parte
dos motores de uma planta sem monitoramento. Nós sensores de baixo custo podem ampliar essa
cobertura, porém os trabalhos que os propõem raramente informam a banda e o piso de ruído do
sensor. Este trabalho investigou quais modos de falha uma cadeia de medição de baixo custo, com
acelerômetro MEMS amostrado a \SI{1}{\kilo\hertz}, microcontrolador ESP32 e enlace LoRa, permite
detectar, e quanto as decisões de estimação alteram o desempenho relatado. O objetivo foi medir a
detecção por modo de falha e calibrar o limiar de um detector Isolation Forest para uma taxa de
falso alarme fixada. Os experimentos usaram \num{1951} gravações de bancada do conjunto público
MAFAULDA, decimadas para \SI{1}{\kilo\hertz}, e sinais sintetizados, com a integração testada
contra um valor analítico. Em dado real, a \SI{2,0 \pm 1,3}{\percent} de falso alarme, a detecção
foi de \SI{75,0 \pm 3,5}{\percent} em desbalanceamento, de \SI{14,8}{\percent} e
\SI{20,7}{\percent} em defeito de rolamento sem massa sobreposta e de \SI{5,0}{\percent} e
\SI{7,2}{\percent} em desalinhamento; com o limiar padrão, o falso alarme foi de \SI{31,2 \pm
4,2}{\percent}. Em sinal sintetizado, a soma acumulada sem remoção da média errou a velocidade
eficaz em \SI{73}{\percent}, e em \SI{716}{\percent} com viés contido na especificação do sensor,
enquanto a integração no domínio da frequência ficou a menos de \SI{0,01}{\percent} do valor
analítico; a área parcial sob a curva ROC caiu até \SI{20,3}{\percent} entre as cadeias de
\SI{20}{\kilo\hertz} e \SI{1}{\kilo\hertz}, contra \SI{3,7}{\percent} da área global; e um
classificador supervisionado atingiu \SI{99,0}{\percent} de acurácia sob partição aleatória de
janelas e \SI{63,2}{\percent} sob partição por ensaio. \emandamento{Em duas montagens de bancada,
sem motor e sem falha induzida, o dispositivo amostrou a cerca de \SI{999,8}{\hertz} em três ciclos,
e o receptor aceitou \num{270} pacotes.} Os resultados indicam que, com esse detector e quatro
descritores, a cadeia permite detectar desbalanceamento sob taxa de falso alarme fixada e detecta
pouco defeito de rolamento sem desbalanceamento sobreposto. Indicam também que o desempenho de
dispositivos dessa classe só pode ser comparado quando a banda, o limiar e o protocolo de validação
são informados.

\vspace{\onelineskip}
\noindent
\textbf{Palavras-chave}: manutenção preditiva; análise de vibração; detecção de anomalias;
sistemas embarcados; instrumentação de baixo custo.
\end{resumo}
